\documentclass[12pt,a4paper]{article}

% --- Packages ---
\usepackage[margin=2.5cm]{geometry}
\usepackage{amsmath}
\usepackage{amssymb}
\usepackage{graphicx}
\usepackage{booktabs}
\usepackage{hyperref}
\usepackage{siunitx}

% --- Title ---
\title{Near-Infrared Photometry of GX~339-4:\\
       {\small A ZOGY Difference Imaging Pipeline}}
\author{Eoin Trainor \\ University College Cork}
\date{\today}

% ============================================================
\begin{document}
\maketitle

% --- Abstract -----------------------------------------------
\begin{abstract}
We present a Ks-band differential photometry pipeline for the black hole
low-mass X-ray binary GX~339-4, using 234 frames across nine observing
blocks obtained with HAWK-I at the ESO VLT over a 67-day baseline in
2025.  ZOGY proper image subtraction isolates the time-variable flux of
the donor star from contaminating field sources in this crowded field.
After per-OB detrending to remove slow accretion variability, the
phase-folded lightcurve shows no significant ellipsoidal modulation above
the frame-to-frame scatter, suggesting the system was in an active or
semi-active state during the campaign.
\end{abstract}

\tableofcontents
\newpage

% ============================================================
\section{Introduction}
\label{sec:intro}

GX~339-4 is one of the most frequently studied black hole low-mass
X-ray binaries (LMXBs) \cite{Markert1973}.  Despite decades of
observation the black hole mass remains poorly constrained, largely
because the system spends much of its time in an active or
semi-active state rather than in the deep quiescence required for
ellipsoidal modelling.

The near-infrared (NIR) Ks band ($\lambda \approx 2.15\,\mu$m) is
preferred for ellipsoidal modelling because the donor star's
thermal emission dominates over the accretion disc contribution
during quiescence.  The ellipsoidal light curve — arising from the
tidal distortion of the Roche-lobe-filling donor — exhibits two
equal maxima per orbit at quadrature ($\phi = 0.25, 0.75$) and two
minima at conjunction ($\phi = 0.0, 0.5$), and its amplitude
constrains the binary mass ratio and inclination.

The crowded nature of the GX~339-4 field makes conventional aperture
photometry unreliable; ZOGY proper image subtraction
\cite{Zackay2016} removes the contribution of blended field stars
while preserving the photometric noise properties of the difference
image, providing a cleaner measurement of differential flux.

% ============================================================
\section{Observations}
\label{sec:obs}

\subsection{Instrument}

Observations were obtained with HAWK-I at the VLT UT4.  HAWK-I
consists of four Hawaii-2RG detectors arranged in a $2 \times 2$
mosaic with a $15''$ cross-shaped gap, giving a total field of view
of $7.5' \times 7.5'$ and a plate scale of $0.106''$/pixel.
GX~339-4 falls on chip\,1 (bottom-left detector).  All data were
taken in the Ks filter with a science DIT of 10\,s and NDIT~=~9,
yielding a per-frame exposure time of 90\,s after jitter-pattern
co-addition.

\subsection{Observing Log}

Nine observing blocks (OBs) were obtained between 2025 May 17 and
2025 July 23, spanning a baseline of 66.8\,days ($\approx 38$ orbital
periods).  Each OB comprised 26--28 jitter frames, for a total of
236 science frames.  Two frames were rejected after per-OB sigma
clipping (one each in OB\,2 and OB\,8), leaving 234 frames in the
final lightcurve.  Table~\ref{tab:obs} summarises each OB.  The
mean seeing over the campaign was $0.71''$ (FWHM), ranging from
$0.42''$ to $1.03''$.

\begin{table}[h]
  \centering
  \caption{Summary of observing blocks.  FWHM is the mean PSF width
           of the science frames for that OB measured in the ZOGY
           pipeline.}
  \label{tab:obs}
  \begin{tabular}{lcccc}
    \toprule
    OB & Date (UT) & Frames & Good & Mean FWHM ($''$) \\
    \midrule
    OB~1 & 2025-05-17 & 26 & 26 & 0.85 \\
    OB~2 & 2025-06-02 & 26 & 25 & 0.82 \\
    OB~3 & 2025-06-04 & 26 & 26 & 0.55 \\
    OB~4 & 2025-06-19 & 26 & 26 & 0.76 \\
    OB~5 & 2025-07-03 & 26 & 26 & 0.70 \\
    OB~6 & 2025-07-14 & 26 & 26 & 0.75 \\
    OB~7 & 2025-07-18 & 26 & 26 & 0.61 \\
    OB~8 & 2025-07-20 & 26 & 25 & 0.83 \\
    OB~9 & 2025-07-23 & 28 & 28 & 0.55 \\
    \bottomrule
  \end{tabular}
\end{table}

% ============================================================
\section{Method: ZOGY Subtraction}
\label{sec:zogy}

\subsection{Calibration and Alignment}

Raw science frames were dark-subtracted using a master dark scaled
to the science DIT, flat-fielded with a Ks master flat, and
bad-pixel corrected.  Frames within each OB were aligned to a common
reference astrometric grid using SExtractor source catalogues and an
affine transformation, then stacked to form a per-OB co-add used
for constructing the ZOGY reference.

\subsection{ZOGY Difference Imaging}

The ZOGY proper difference image $\hat{D}$ is given by
\cite{Zackay2016}:

\begin{equation}
  \hat{D} = \frac{F_r\,\hat{P}_r\,\hat{N} - F_n\,\hat{P}_n\,\hat{R}}
                 {\sqrt{F_r^2\,\sigma_n^2\,|\hat{P}_r|^2 +
                        F_n^2\,\sigma_r^2\,|\hat{P}_n|^2}}
  \label{eq:D}
\end{equation}

where $R$ and $N$ are the reference and science images, $P_r$,
$P_n$ their PSFs, $\sigma_r$, $\sigma_n$ their background RMS
values, and $F_r$, $F_n$ their photometric flux normalisations.
PSFs were modelled as circular Gaussians fitted from bright
isolated stars in each frame.

The matched-filter detection statistic $S$ is computed from
$\hat{D}$ by convolving with the science PSF; its value at the
target position provides a SNR-optimal estimate of the differential
flux per frame.

\subsection{Differential Photometry}

For each science frame, $S_\mathrm{target}$ was extracted at the
centroid of GX~339-4 and $S_\mathrm{ref,i}$ at the positions of
eight non-variable reference stars spanning Ks\,$= 14$--$17$\,mag.
Frames were sigma-clipped per OB at $3\sigma$ to remove cosmic-ray
hits and tracking failures.  A per-OB mean was then subtracted from
$S_\mathrm{target}$ to remove OB-to-OB flux offsets caused by
accretion variability and photometric zero-point drift.  The
detrended signal was normalised by its global standard deviation.
Each OB spans $\approx 52$\,min, which is much less than the
$P = 1.7587$\,d orbital period, so the intra-OB ellipsoidal slope
is negligible and this detrending does not suppress the ellipsoidal
signal.

% ============================================================
\section{Results}
\label{sec:results}

Figure~\ref{fig:lc} shows the phase-folded differential lightcurve
of GX~339-4, folded on the Heida+2017 ephemeris
($T_0 = 57529.397$\,MJD, $P = 1.7587$\,d).  The 234 good frames
are distributed over orbital phases $\phi \approx 0.28$--$0.94$;
the region around inferior conjunction ($\phi \approx 0$) was not
sampled during this campaign.

The binned means (20 phase bins, minimum 3 points per bin) are
consistent with zero within their uncertainties across all sampled
phases.  No significant ellipsoidal modulation is detected: the
expected double-humped signal with maxima at quadrature and minima
at conjunction ($\Delta m \sim 0.1$--$0.2$\,mag in quiescence) is
not present above the frame-to-frame scatter.  The scatter in
individual $S_\mathrm{target}$ values is dominated by the large
OB-to-OB flux offset (raw $\sigma \approx 8 \times 10^8$ matched-filter
units), consistent with the system being in a variable
semi-active state rather than deep quiescence during the 2025
campaign.  The eight reference stars show comparable scatter in $S$,
confirming that the differential noise is driven by PSF-matching
systematics and residual sky variability rather than astrophysical
variability in GX~339-4 alone.

\begin{figure}[h]
  \centering
  \includegraphics[width=\textwidth]{C:/Astronomy/GX 339-4/GX 339-4 Output/difference/lightcurve/lc05_phase_folded_paper.png}
  \caption{Phase-folded Ks-band differential lightcurve of GX~339-4
           on the Heida+2017 ephemeris ($P = 1.7587$\,d,
           $T_0 = 57529.397$\,MJD).  Individual frames are shown as
           coloured points (one colour per OB); black points with
           error bars are the binned mean $\pm$ standard error in 20
           equal phase bins.  Blue dotted lines mark inferior and
           superior conjunction ($\phi = 0, 0.5$); red dotted lines
           mark quadrature ($\phi = 0.25, 0.75$).}
  \label{fig:lc}
\end{figure}

% ============================================================
\section{Conclusions}
\label{sec:conclusions}

We have demonstrated a Ks-band ZOGY difference imaging pipeline for
GX~339-4 using nine HAWK-I observing blocks spanning 67\,days in
2025.  The pipeline reliably extracts per-frame differential
photometry from a crowded NIR field, with a mean seeing of
$0.71''$ and 234 frames surviving quality cuts.

No significant ellipsoidal modulation is detected in the
phase-folded lightcurve.  The dominant signal is OB-to-OB flux
variability at the $\sim 10^8$--$10^9$ matched-filter unit level,
consistent with the system remaining in an active or semi-active
accretion state throughout the campaign.  Additionally, the orbital
phase coverage does not include the inferior conjunction region
($\phi \lesssim 0.28$), which is required for a complete ellipsoidal
fit.

Future work should target GX~339-4 during a confirmed X-ray
quiescent epoch and obtain full orbital phase coverage, ideally
with OBs distributed to sample all phases including conjunction.
With the pipeline validated here, such a dataset could yield a
robust ellipsoidal solution and place new constraints on the binary
inclination and mass ratio.

% ============================================================
\begin{thebibliography}{9}

\bibitem{Markert1973}
  Markert T.~H. et al., 1973, ApJ, 184, L67.

\bibitem{Zackay2016}
  Zackay B., Ofek E.~O., Gal-Yam A., 2016, ApJ, 830, 27.

\end{thebibliography}

\end{document}
